\section{Problem Definition \& Notations}
\label{sec:hybrid-prelims}
% \subsection{Problem Definition \& Notations}
\paragraph{Input} 
 Our dataset is a collection of human subject's features and their phenotypic outcomes, which is represented by the pair $(X, Y)$ for each subject. $X \in \mathbb{R}^{N \times d}$ represents the features for each subject, where $N$ is the number of brain regions and $d$ is the feature size. Consistent with previous works \citep{kan2022brain,li2021braingnn}, the features are Pearson correlations derived from fMRI time series. $Y \in \mathbb{R}$ denotes the phenotypic outcome, such as the intelligent quotient. Full implementation details and hyperparameter settings are provided in \cite{qiu2023learning}.
 
 %For how to obtain features of regions from raw data, see for more details.

\paragraph{Goal}
Based on the input $X$, \hybrid{} aims to learn a weighted hypergraph of the brain, where regions are nodes. To achieve this, \hybrid{} identifies a collection of hyperedges $H = (\boldsymbol{h}^1, \boldsymbol{h}^2, \cdots, \boldsymbol{h}^K)$, and assigns weights $\boldsymbol{w} = [w^1, w^2, \cdots, w^K]^T$ for all hyperedges. These hyperedges and their weights, which represent strengths of hyperedges, are expected to be the most informative towards $Y$ yet the least redundant. 

\paragraph{Representation of Hyperedges}
As mentioned before, we use $H$ to denote the collection of hyperedge structures and $\boldsymbol{h}^k$ to denote the $k$-th hyperedge. To associate hyperedges with node memberships, we use the following representation for a hyperedge:
\begin{equation}
    \boldsymbol{h}^k = \boldsymbol{m}^k \odot X \in \mathbb{R}^{N\times d},
\end{equation}
where $\boldsymbol{m}^k \in \{0,1\}^N$ is a mask vector and $\odot$ denotes broadcasting element-wise multiplication. In other words, each $\boldsymbol{h}^k$ is a randomly row-zeroed version of $X$.